\documentclass[../main.tex]{subfiles}

\begin{document}

\section{Parameteros de control}

\subsection{Red polimérica}

\paragraph{Partículas patchy}
Controla la definición de los monomeros a usar para representar/construir una red polimérica.

\begin{itemize}
    \item \(D_c\): Diámetro de la partícula central
    \item \(D_p\): Diámetro de la partícula patchy
    \item \(v\): Cantidad de partículas patchy (valencia)
    \item \(s\): Separación del centro de la partícula patchy con el centro de la partícula central
    \item \(\alpha\): Angulo entre los centros de de las partículas patchy. 
\end{itemize}

\paragraph{Potenciales de interacción}
Controla parte de la interacción entre los monómeros.

\begin{itemize}
    \item \(\epsilon\): Pozo de potencial
    \item \(r_{\min}\): Distancia en donde \(U(r_{\min})=-\epsilon\) 
    \item \(r_c\): Distancia de corte
\end{itemize}


\paragraph{Red polimérica (Hidrogel)}\footnote{Restricción del sistema: \[\phi\leq 0.56\] \[\sum_{v=2}^{N}\chi_v=1\]}
Controlan la topología de la red polimérica.

\begin{itemize}
    \item \(\phi\): Factor de empaquetamiento
    \item \(\chi_v\): Concentración de partículas patchy con valencia \(v\)
\end{itemize}

\subsection{Sistema}

\paragraph{Movimiento Browniano}
Controla la interacción entre la red polimérica con un baño térmico.

\begin{itemize}
    \item \(T\): Temperatura 
    \item \(\mathrm{damp}\): Relacionado a la viscosidad
    \item \(m\): Masa de las partículas
\end{itemize}

\paragraph{Solución computacional}
Controla la estabilidad del método numérico.

\begin{itemize}
    \item \(\Delta t\): Paso temporal 
\end{itemize}

\subsection{Parametros del experimento}

\paragraph{Definición}
Controla el \emph{tipo} de experimento.

\begin{itemize}
    \item \(t_{\mathrm{temp}}\): Tiempo en el que se eleva la temperatura de \(0\) a \(T\).
    \item \(t_{\mathrm{isot}}\): Tiempo en que se deja interactuar las partículas patchy a una misma temperatura
    \item \(N_p\): Número de partículas patchy totales.
\end{itemize}



\section{Análisis del sistema}

Para cada experimento se piensa hacer una gráfica de:
\begin{itemize}
    \item Energía total del sistema con respecto al tiempo
    \item Evolución temporal del factor de estructura 
    \item Conectividad de la red?
    \item Evolución temporal de la cantidad de partículas centrales en el cluster más grande
    \item Velocidad promedio de las partículas centrales con respecto al tiempo?
    \item Cantidad de cadenas sueltas/finales?
    \item Cantidad de ciclos
    \item Distribución de la cantidad de monomeros de valencia 2 enlazados a los monómeros de valencia 4
    \item Direccionalidad de la red
\end{itemize}


\section{Parámetros a variar}

\begin{enumerate}
    \item Temperatura \(T\)
    \item Porcentaje de CrossLinkers \(\chi_4\)
    \item Fracción de empaquetamiento \(\phi\)
    \item Pozo de potencial y distancia de interacción \(\epsilon, r_c\)
    \item Cantidad de patches en monomeros \(v\)
    \item Ubicación de los patches de los monomeros \(\alpha\)
\end{enumerate}


\section{Objetivos}

\begin{itemize}
    \item Caracterizar la red polimérica.
\end{itemize}




\end{document}
